%% ma: L10-B02-0024
%% lop: 10
%% chuong: 1
%% bai: 2
%% dang: 10.2.5
%% loai: TLN
%% muc: TH
%% dap_an: "-"
%% nguon: "Ảnh giáo viên gửi 10/10/2026 (ThS. Nguyễn Hoàng Việt, Chương 2 Tập hợp), A-Đề 1, Phần 3 Câu 20"
%% kien_thuc: "Bài 2 (tập hợp và các phép toán trên tập hợp)"
%% trang_thai: cach-ly
%% ly_do: "Đề yêu cầu hai số (số phần tử của A giao B là 1 và của A hiệu B là 4) nhưng câu trả lời ngắn chỉ nhận một đáp số. Gợi ý: sửa thành Tính tổng số phần tử của A giao B và A hiệu B (đáp số 5), hoặc tách thành hai câu"
\begin{ex}
Cho hai tập hợp $A=\{n\in\mathbb N\mid-2<n\le4\}$ và $B=\{x\in\mathbb Z\mid2x^3+x^2-x=0\}$. Hãy viết số phần tử của tập hợp $A\cap B$, $A\setminus B$.
\shortans{0}
\loigiai{$A=\{0;1;2;3;4\}$. $2x^3+x^2-x=x(2x-1)(x+1)=0\Leftrightarrow x\in\left\{0;\dfrac12;-1\right\}$, nên $B=\{-1;0\}$. $A\cap B=\{0\}$ có $1$ phần tử; $A\setminus B=\{1;2;3;4\}$ có $4$ phần tử.}
\end{ex}
